\paragraph{The odd-prime-power cyclic-support lemma.}
We use the following extension of the odd-prime argument in
\cite[Theorem~13 and its proof, pp.~14--19]{RootSchmidtWingberg1998}.
Let \(p\) be odd, let \(m=p^r\), and let
\(\Omega/K/k\) be finite extensions of number fields with both
\(\Omega/k\) and \(K/k\) Galois.  Suppose \(\mu_m\subset K\).
Let \(A\) be a finite \(\mathbf Z/m[\operatorname{Gal}(K/k)]\)-module,
and let \(T\) contain the primes ramified in \(\Omega/k\),
the primes above \(p\), and the infinite places.  Put
\[
 S=\operatorname{cs}(\Omega/k)\cup T.
 \tag{A.130.S1}
\]
Here \(k_S\) is the maximal extension unramified outside \(S\);
all cohomology is continuous.  Suppose
\(y\in H^1(k_S/K,A)\) is unramified at every place of \(T(K)\)
and is zero at the places above
\(\operatorname{Ram}(K/k)\cup S_p\cup S_\infty\).
There exists \(z\in H^1(k_S/K,A)\) with the following exact
properties:
\[
\begin{aligned}
 z_{\mathfrak p}&=y_{\mathfrak p}
                         &&(\mathfrak p\in T(K)),\\
 z_{\mathfrak p}\text{ ramified},\ \mathfrak p\notin T(K)
 &\ \Longrightarrow\
 \begin{cases}
 z_{\mathfrak p}\text{ is split by a cyclic extension of }K_{\mathfrak p},\\
 z_{\sigma\mathfrak p}=0
       \quad\text{for all }1\ne\sigma\in\operatorname{Gal}(K/k).
 \end{cases}
\end{aligned}
 \tag{A.130.S2}
\]
Consequently \(x=\operatorname{cor}_{K/k}z\) satisfies
\[
 x_v=(\operatorname{cor}_{K/k}y)_v\quad(v\in T),
 \qquad x_v\text{ is cyclic}\quad(v\notin T).
 \tag{A.130.S3}
\]
The cyclicity assertion means that the actual local cohomology
class becomes zero over one finite cyclic extension.
In particular, at a place splitting \(K\), it means that the
image of its local homomorphism is a cyclic subgroup of \(A\).
The construction gives a stronger assertion when the underlying
group \(A\) is free over \(\mathbf Z/m\): at every new ramified
place the inertia image has exact order \(m\) and generates
the entire local image.  In tame coordinates \(t,c\in A\)
this gives the literal relation \(c=zt\) for some
\(z\in\mathbf Z/m\), with \(t\) of order \(m\).

\emph{The finite duality input.}
We spell out the part of the proof where a change of coefficient
exponent could otherwise be hidden.  For a finite Galois field
\(\Omega_0/K/k\), a set \(T_0\) as above, and
\(S_0=\operatorname{cs}(\Omega_0/k)\cup T_0\), the image of
\[
 H^1(k_{S_0}/K,\mu_e)\longrightarrow
       \prod_{\mathfrak p\in T_0(K)}H^1(K_{\mathfrak p},\mu_e),
                       \qquad e\mid m,
 \tag{A.130.S4}
\]
is the annihilator, under the sum of local Tate pairings, of
\(\operatorname{Hom}(\operatorname{Gal}(\Omega_0/K),\mathbf Z/e)\).
This is the finite-coefficient Poitou--Tate sequence with
unramified conditions outside \(T_0\), followed by Chebotarev.
For clarity, in a finite enlargement of \(T_0\) the annihilating
dual classes are unramified outside \(T_0\), because their
local restrictions at every added place are zero.  They
therefore lie in the one finite group
\(H^1(k_{T_0}/K,\mathbf Z/e)\).
Intersecting their kernels as added primes run through
\(\operatorname{cs}(\Omega_0/k)\) stabilizes in this finite group.
Chebotarev identifies the stabilized subgroup with exactly
the characters factoring through \(\Omega_0/K\): any character
not factoring through that field is nonzero on some Frobenius
whose underlying \(k\)-prime splits completely in \(\Omega_0\).
This proves (A.130.S4), including actual membership in the
image rather than only membership in its closure.

We also use its ramification-coordinate version:
prescribed full classes at \(T_0(K)\), and prescribed classes
modulo the unramified subgroup at finitely many other places,
are realized precisely when they annihilate
\(H^1(k_{T_0}/K,\mathbf Z/e)\).
Its pairing at a new unramified prime sends a valuation-one
Kummer class to evaluation of the dual character on Frobenius.
This is the same finite Poitou--Tate sequence and is the upper
exact row used on page~15 of the cited proof.

\emph{A scalar class of odd prime-power order.}
First take \(A=\mu_e\), where \(e\mid m\), viewed over \(K\).
Fix the usual compatible local Hilbert pairings with values
in \(\mu_e\); choose their order so that pairing an
unramified Kummer class with a valuation-one class evaluates
the former on Frobenius.
Their global product is one, and the unramified subgroup
is isotropic.  We write cohomology additively.

We construct classes \(z_i\in H^1(k_S/K,\mu_e)\) and primes
\(\mathfrak p_i\), with distinct underlying \(k\)-primes
outside \(T\), such that
\[
\begin{gathered}
 z_i\text{ has valuation }1\bmod e\text{ at }\mathfrak p_i,\\
 z_i\text{ is unramified at every other place};\\
 (z_i)_{\mathfrak p}=\tfrac12y_{\mathfrak p}
                                     \quad(\mathfrak p\in T(K));\\
 (z_j)_{\mathfrak p_i}=0,\qquad
 (z_j)_{\sigma\mathfrak p_i}=-(z_i)_{\sigma\mathfrak p_i},\\
 j>i,\qquad 1\ne\sigma\in\operatorname{Gal}(K/k).
\end{gathered}
 \tag{A.130.S5}
\]
The inverse of two exists modulo \(e\).
Every underlying prime is chosen to split completely in
\(\Omega\), so its different \(\operatorname{Gal}(K/k)\)-conjugates
are different places of \(K\).

Here is the precise induction.  Add all the earlier prime
orbits to \(T\), obtaining \(T_i\), and prescribe the last
two lines of (A.130.S5), with zero at each earlier distinguished
\(\mathfrak p_j\).  These are all unramified classes.
Their duality obstruction is a functional on the finite group
\(H^1(k_{T_i}/K,\mathbf Z/e)\).
Its restriction to the characters of \(\Omega/K\) is zero.
Indeed the earlier prime orbits split in \(\Omega\), while
at \(T(K)\) the functional is one half of that of the actual
global class \(y\); reciprocity for \(y\) has no contribution
outside \(T(K)\), because there \(y\) is unramified or the
prime splits in \(\Omega\).

Let \(F_i/K\) be the finite abelian extension cut out by
\(H^1(k_{T_i}/K,\mathbf Z/e)\).
The negative of the obstruction is an element of
\(\operatorname{Gal}(F_i/K)\) fixing \(F_i\cap\Omega\).
It therefore lifts to an element of
\(\operatorname{Gal}(F_i\Omega/K)\) acting identically on
\(\Omega\).  Chebotarev chooses a new underlying prime
splitting \(\Omega/k\) and a place \(\mathfrak p_{i+1}\)
above it with this Frobenius.  One may exclude all earlier
primes and any additional finite set.
Prescribe valuation one at that place and zero ramification
at its other conjugates.  Its Frobenius functional cancels
the obstruction.  The ramification-coordinate exact sequence
then supplies \(z_{i+1}\), proving (A.130.S5).

For each \(i\) record the finite vector
\[
 \lambda_i=\bigl(
 z_i(\operatorname{Frob}_{\sigma\mathfrak p_i})
                \bigr)_{1\ne\sigma\in\operatorname{Gal}(K/k)}
        \in\prod_{\sigma\ne1}\mu_e.
 \tag{A.130.S6}
\]
The displayed local restrictions are unramified, so these
values are defined.  Choose \(i<j\) with
\(\lambda_i=\lambda_j\), and put \(z=z_i+z_j\).
At \(T(K)\) this is \(y\).  At the conjugates of
\(\mathfrak p_i\), other than \(\mathfrak p_i\) itself,
the last line of (A.130.S5) gives zero directly.

For completeness the signs at the other distinguished
prime follow from two literal Hilbert products.
Fix \(\sigma\ne1\), and write
\(u=z_i(\operatorname{Frob}_{\sigma\mathfrak p_i})\).
In the product for \((z_i,\sigma z_i)\), only
\(\mathfrak p_i\) and \(\sigma\mathfrak p_i\) can contribute:
at \(T(K)\) both arguments are unramified, and elsewhere
both are unramified except at those two places.
The value at \(\sigma\mathfrak p_i\) is \(u\), so the value
at \(\mathfrak p_i\) is \(u^{-1}\).
By (A.130.S5) at \(\sigma^{-1}\mathfrak p_i\),
\[
       (\sigma z_j)_{\mathfrak p_i}
                         =-(\sigma z_i)_{\mathfrak p_i}.
 \tag{A.130.S7}
\]
The product for \((z_i,\sigma z_j)\) consequently has value
\(u\) at \(\mathfrak p_i\), and value \(u^{-1}\) at
\(\sigma\mathfrak p_j\).
The latter is exactly
\(z_i(\operatorname{Frob}_{\sigma\mathfrak p_j})\).
Equality of the vectors (A.130.S6) says that
\(z_j(\operatorname{Frob}_{\sigma\mathfrak p_j})=u\).
Their product is one.  Thus
\(z_{\sigma\mathfrak p_j}=0\) for every \(\sigma\ne1\).
At each distinguished prime the ramification coordinate
of \(z\) is still one modulo \(e\).
The scalar local homomorphism has cyclic target \(\mu_e\),
so it is cyclic.  This proves (A.130.S2) in the scalar case.
The argument used finiteness, Hilbert reciprocity, and the
inverse of two; it did not use that \(\mathbf Z/e\) is a field.

\emph{The bounded-character extension in a compositum.}
The following exact group fact supplies the module induction.
Let \(G_1\) be finite, let \(B\) be finite abelian killed by \(e\),
and let
\[
 G=G_1\times_I B,\qquad
       \pi_1:G\longrightarrow G_1,\quad\pi_2:G\longrightarrow B
 \tag{A.130.S8}
\]
be a fibre product for surjections onto a common quotient \(I\).
Every character \(\chi:G\to\mathbf Z/e\) has an expression
\[
                 \chi=\chi_1\circ\pi_1+\chi_2\circ\pi_2.
 \tag{A.130.S9}
\]
To prove this, identify \(\ker\pi_1\) with the subgroup
\(\ker(B\to I)\) of \(B\).
Its character \(\chi|_{\ker\pi_1}\) extends to a character
of \(B\) with values in \(\Q/\Z\).
This finite-character extension follows, for example, by
adjoining generators one at a time and taking the required
roots in the divisible group \(\Q/\Z\).
Since \(eB=0\), the image of every such extension is killed
by \(e\).  It is therefore a character \(\chi_2:B\to\mathbf Z/e\).
Subtracting its pullback from \(\chi\) makes the restriction
to \(\ker\pi_1\) zero; the difference descends to \(\chi_1\).
This proves (A.130.S9) without any linear-disjointness assertion.

\emph{Induction on the underlying cyclic summands.}
The construction of \(z\) and property (A.130.S2) concern
only the constant coefficient group over \(K\), its local
images, and the permutation of the places of \(K\).
We prove this construction for every finite abelian
\(p\)-group over \(K\), without prescribing an action of
\(\operatorname{Gal}(K/k)\) on that group.
The original arithmetic action is used only in the
corestriction at the end.
Decompose the underlying abelian group of \(A\) as a direct
sum of cyclic \(p\)-groups, ordered by nondecreasing exponent.
This is a decomposition over \(K\), where the action is trivial;
no summand is asserted to be stable under
\(\operatorname{Gal}(K/k)\).
Whenever a scalar construction uses a Galois action on
one cyclic summand, transport the cyclotomic action
through its chosen identification with the appropriate
\(\mu_{p^e}\).  Its conclusion about zero restrictions
at conjugate places does not depend on this temporary
extension of the trivial \(G_K\)-action.
Write the next induction step as
\(A=A'\oplus\mu_e\), using a chosen isomorphism of the last
cyclic group with \(\mu_e\).  The exponent of \(A'\) divides \(e\).
Apply the induction assertion over \(K\) to \(y'\) and obtain
\(z'\).  Let \(K'/K\) be the finite abelian extension cut out
by this homomorphism, and let \(\widetilde K/k\) be its
normal closure.  The group
\(\operatorname{Gal}(\widetilde K/K)\) embeds in a product
of the conjugate groups \(\operatorname{Gal}(K'/K)\).
It is therefore abelian of exponent dividing \(e\).
It is unramified at \(T(K)\), and any further ramification
lies above primes splitting in \(\Omega/k\).

Put
\[
 \Omega'=\Omega\widetilde K,\qquad
 T'=T\cup\operatorname{Ram}(\widetilde K/k),\qquad
 S'=\operatorname{cs}(\Omega'/k)\cup T'.
 \tag{A.130.S10}
\]
In particular \(S'\subset S\).
Prescribe a scalar local class \(\eta\) equal to \(y''\)
on \(T(K)\) and zero on \(T'(K)\setminus T(K)\).
Every \(\mathbf Z/e\)-valued character of
\(\operatorname{Gal}(\Omega'/K)\) is a sum of a character
from \(\Omega/K\) and one from \(\widetilde K/K\), by
(A.130.S8)--(A.130.S9).
The first character pairs to zero with \(\eta\) by global
reciprocity for the original global component \(y''\):
outside \(T(K)\), it vanishes at split primes and pairs
trivially with unramified classes elsewhere.
The second character is unramified at \(T(K)\), where
\(\eta\) is unramified, and \(\eta\) is zero on the added places.
It too pairs to zero.
The exact image statement (A.130.S4) therefore gives
\[
       \widetilde y\in H^1(k_{S'}/K,\mu_e),
                 \qquad \widetilde y|_{T'(K)}=\eta .
 \tag{A.130.S11}
\]

Apply the scalar construction to \(\Omega'/K/k\) and
\(\widetilde y\), obtaining \(z''\).
Set \(z=z'+z''\), viewed in the actual direct sum \(A\).
At \(T(K)\) it equals \(y\).
At a ramified place of \(z'\) outside \(T(K)\), the scalar
component is zero, because that place belongs to
\(T'(K)\setminus T(K)\).
At a new ramified place of \(z''\), its underlying prime
splits in \(\Omega'\), hence in the normal closure of \(K'\).
The entire component \(z'\), at that place and all its
\(\operatorname{Gal}(K/k)\)-conjugates, is zero.
Thus at every new ramified place only one of the two
components is present, and its nontrivial conjugate
restrictions vanish.  This proves (A.130.S2) inductively.
If every cyclic summand has order \(m\), its scalar
construction has inertia coordinate one modulo \(m\)
at both distinguished primes.  At a new prime belonging
to that summand, all earlier components vanish because
the prime splits their normal closure, and all later
components are prescribed to vanish there.  The inertia
therefore has order \(m\) and generates the full local
cyclic image.

Finally use the actual corestriction
\(H^1(k_S/K,A)\to H^1(k_S/k,A)\), with the originally
given \(\operatorname{Gal}(K/k)\)-action on \(A\).
At \(T\) its local Mackey formula gives the first equality
in (A.130.S3).
Outside \(T\), a ramified prime splits in \(\Omega\) and
hence in \(K\).  The same formula is a sum over its
distinct \(K\)-places.  Property (A.130.S2) leaves only
one nonzero ramified term, whose image is cyclic; its
conjugate by the actual coefficient action still has
cyclic image.  In the free case the coefficient action
is an automorphism, so it preserves the asserted order
\(m\) of inertia and the equality of the inertia and
decomposition images.  If no ramification remains, the class is
unramified and is killed by a finite unramified cyclic
extension.  This proves (A.130.S3) and completes the
odd-prime-power extension.

